\section{Proporción}
\subsection{Teoría}
Se define a la razón de proporcionalidad como lo siguiente:
\begin{definicion}
    Si a y b son dos cantidades directamente proporcionales, la razón $\frac{a}{b}$ recibe el nombre de razón de proporcionalidad, la cual siempre es constante.
\end{definicion}
\begin{ejemplo}
    Si 18 libros de ciencia cuestan \$1260, la razón de proporcionalidad es de 70, ya que
    \begin{equation*}
        \frac{1260}{18} = 70
    \end{equation*}
\end{ejemplo}
La proporción se define como la igualdad entre dos razones, de tal manera que si a y b tienen una razón r, c y d tienen una razón r entonces:
\begin{equation*}
    \frac{a}{b}=\frac{c}{d}
\end{equation*}
\begin{ejemplo}
    3 es a 6 como 8 es a 16, se escribe
    \begin{equation*}
        \frac{3}{6} = \frac{8}{16}
    \end{equation*}
    al simplificar cada fracción se obtiene $\frac{1}{2}$
\end{ejemplo}